%% kapa: exam=1 round=30 year=2019 period=1 form=A subject=economics subject_ko=경제학원론 seq=14 number=54
%% source: Q-Net 감정평가사 기출문제 (cst00309) artlSeq=5205205 file=extracted/1008/5205205/2019년 제30회 감정평가사 1차시험 문제지/2019년 제30회 감정평가사 1차시험 문제지/2019년도 제30회 감정평가사 1차 1교시 A형.hwp
%% conversion: kordoc 4.24.51 -> tools/gen1.py (auto); review=proofread-sonnet; flags=none
\begin{question}{54}
임금의 보상격차(compensating differential)에 관한 설명으로 옳지 않은 것은?
\begin{choices}
  \item 근무조건이 좋지 않은 곳으로 전출되면 임금이 상승한다.
  \item 성별 임금 격차도 일종의 보상격차이다.
  \item 비금전적 측면에서 매력적인 일자리는 임금이 상대적으로 낮다.
  \item 물가가 높은 곳에서 근무하면 임금이 상승한다.
  \item 더 높은 비용이 소요되는 훈련을 요구하는 직종의 임금이 상대적으로 높다.
\end{choices}
\end{question}
\answer{2}
